% some parameters for customization
\def\shadowshift{2pt,-5pt}
\def\shadowradius{10pt}

% this draws a shadow under a rectangle node
\newcommand\drawshadow[2]{
    \begin{pgfonlayer}{shadow}
        \shade[white,inner color=#2,outer color=white] ($(#1.south west)+(\shadowshift)+(\shadowradius/2,\shadowradius/2)$) circle (\shadowradius);
        \shade[white,inner color=#2,outer color=white] ($(#1.north west)+(\shadowshift)+(\shadowradius/2,-\shadowradius/2)$) circle (\shadowradius);
        \shade[white,inner color=#2,outer color=white] ($(#1.south east)+(\shadowshift)+(-\shadowradius/2,\shadowradius/2)$) circle (\shadowradius);
        \shade[white,inner color=#2,outer color=white] ($(#1.north east)+(\shadowshift)+(-\shadowradius/2,-\shadowradius/2)$) circle (\shadowradius);
        \shade[top color=#2,bottom color=white] ($(#1.south west)+(\shadowshift)+(\shadowradius/2,-\shadowradius/2)$) rectangle ($(#1.south east)+(\shadowshift)+(-\shadowradius/2,\shadowradius/2)$);
        \shade[left color=#2,right color=white] ($(#1.south east)+(\shadowshift)+(-\shadowradius/2,\shadowradius/2)$) rectangle ($(#1.north east)+(\shadowshift)+(\shadowradius/2,-\shadowradius/2)$);
        \shade[bottom color=#2,top color=white] ($(#1.north west)+(\shadowshift)+(\shadowradius/2,-\shadowradius/2)$) rectangle ($(#1.north east)+(\shadowshift)+(-\shadowradius/2,\shadowradius/2)$);
        \shade[white,right color=#2,left color=white] ($(#1.south west)+(\shadowshift)+(-\shadowradius/2,\shadowradius/2)$) rectangle ($(#1.north west)+(\shadowshift)+(\shadowradius/2,-\shadowradius/2)$);
        \filldraw[fill=#2, draw=#2] ($(#1.south west)+(\shadowshift)+(\shadowradius/2,\shadowradius/2)$) rectangle ($(#1.north east)+(\shadowshift)-(\shadowradius/2,\shadowradius/2)$);
    \end{pgfonlayer}
}

% create a shadow layer, so that we don't need to worry about overdrawing other things
\pgfdeclarelayer{shadow} 
\pgfsetlayers{shadow,main}

\newcommand{\linenumber}[1]{\texttt{\textcolor{gray}{#1:\ \ }}}
\newcommand{\texttug}[1]{\texttt{\textcolor{tug}{#1}}}
\newcommand{\texttugb}[1]{\texttt{\textcolor{tugb!120}{#1}}}
\newcommand{\textblack}[1]{\texttt{#1}}
\newcommand{\textorang}[1]{\texttt{\textcolor{color4}{#1}}}
\newcommand{\textgreen}[1]{\texttt{\textcolor{tugb!120}{#1}}}
    
\begin{tikzpicture}[
    node distance=-0.1mm,
    %text width=0.65\columnwidth,
    textstyle/.style={font=\ttfamily, text width=0.85\columnwidth},
    every node part/.style={font=\texttt},] 
    
\node[rotate=90,anchor=south, align=center, minimum width=193pt, text width=193pt, minimum height=15pt, fill=lightgray, draw=lightgray](sys_msg){System Rules};

\node[rotate=90,anchor=south, align=center, minimum width=128pt, text width=128pt, minimum height=15pt, fill=lightgray!40, left=of sys_msg, , draw=lightgray!40](usr_msg){Catalog Data};

\node[textstyle, right=of sys_msg, xshift=8pt, yshift=-9pt, draw=lightgreencode,text width=0.49\columnwidth, inner ysep=2.2mm, very thick] {};
\node[textstyle, right=of sys_msg, xshift=7pt, yshift=-9pt, anchor=west] (ruleTask){\#\#\# \textbf{Task:} Generate a data science pipeline in Python 3.10.};

\node[textstyle, anchor=west, below=of ruleTask] (rule){\#\#\# \textbf{Input:} A dataset in CSV format, a schema that describes the columns and data types of the dataset, and a data profiling info that summarizes the statistics and quality of the dataset.};
    
\node[textstyle,anchor=west, below=of rule] (rule){\#\#\# \textbf{Output:} A pipeline code that performs the following steps:};

\node[textstyle, below=of rule.south west, anchor=north west, draw=lightgreencode, fill=lightgreencode!30,minimum width=0.4\columnwidth, text width=0.4\columnwidth, inner ysep=2mm, thick] (himport) {};
\node[textstyle,anchor=west, below=of rule] (ruleImport){\ \textbf{\#1:} Import the necessary libraries and modules.};

\node[textstyle,  below=of ruleImport.south west, anchor=north west, yshift=-1pt, draw=lightgreencode,fill=lightgreencode!20, inner ysep=2mm, thick, text width=0.72\columnwidth] (hdatasets){\ \\ \ };
\node[textstyle, below=of ruleImport] (rule){\ \textbf{\#2:} Load the training and test datasets. For the training data, utilize the variable "train\_data.csv", and for the test data, employ the variable "test\_data.csv".};
\node[textstyle, below=of rule] (rule){\ \textbf{\#3:} The user will provide the Schema, and Data Profiling Info of the dataset with columns appropriately named as attributes.}; 
\node[textstyle, below=of rule.south west, anchor=north west, yshift=0pt, draw=lightgreencode, fill=lightgreencode!30,text width=0.6\columnwidth, inner ysep=2mm,thick] (hclean) {};  
\node[textstyle, below=of rule] (rule){\ \textbf{\#4:} Perform missing value imputation for features 'Address' and 'Zip'.};

\node[textstyle, below=of rule.south west, anchor=north west, yshift=-1.5pt, draw=lightgreencode, fill=lightgreencode!30,text width=.65\columnwidth, inner ysep=2.5mm,thick] (hfe){\\ \ };
\node[textstyle, below=of rule] (rule){\ \textbf{\#5:} Perform feature extraction (dataset contains categorical values).};
\node[textstyle, below=of rule] (rule){\ \textbf{\#6:} Perform feature selection (redundant columns e.g., 'Address' and 'Zip').};
\node[textstyle, below=of rule.south west, anchor=north west, yshift=-1pt, draw=lightgreencode, fill=lightgreencode!30,text width=0.5\columnwidth, inner ysep=2.5mm,thick] (hmodel) {\\ \ };
\node[textstyle, below=of rule] (rule){\ \textbf{\#7:} Select an appropriate ML algorithm.};
\node[textstyle, below=of rule] (rule){\ \textbf{\#8:} Assign a default value to a particular hyperparameter.};   
\node[textstyle, below=of rule.south west, anchor=north west, draw=lightgreencode, fill=lightgreencode!30, minimum width=120pt, text width=120pt] (rule-ev){\ \textbf{\#9:} Evaluate the model.};
    
\draw[densely dotted, very thick, gray] ($(rule-ev.south west)+(-14pt,-1pt)$) -- ($(rule-ev.south east)+(+350pt,-1pt)$);
    
\node[textstyle, below=of rule-ev.south west, anchor=north west, yshift=-5pt] (rule){\textbf{\#\#\#  Dataset Description:} A The dataset was obtained from multiple sources, including surveys, job posting sites, and other publicly available sources.A total of 6704 data points were collected.The dataset included fivevariables: age, experience, job role, and education level and salary.};
\node[textstyle, below=of rule] (rule){\textbf{\#\#\# Schema, and Data Profiling Info:}};
\node[textstyle, below=of rule, xshift=2pt, draw=tugb, fill=tugb!10, inner ysep=10.7mm,thick] (hschema) { };

\node[textstyle, below=of rule] (rule){ \shortstack[l]{
\# Experience (string), distinct [ 60\% ], missing [ 0\% ], categorical-vals [1 year, 2 years, 3 years] \vspace{0.3mm}\\ 
\# Gender (string), distinct [ 40\% ], missing [ 0\% ], categorical-vals [ Male, Female ] \vspace{0.3mm}\\ 
\# Address (string), distinct [ 40\% ], missing [ 20\% ], categorical-vals [ CA, TX ] \vspace{0.3mm} \\
\# Zip (string), distinct [ 40\% ], missing [ 40\% ], categorical-vals [ 7050, 7871 ] \vspace{0.7mm} \\
\\ \\ \
}};

\node[textstyle, below=of rule.south west, anchor=north west, inner xsep=0.1mm, xshift=2pt ,yshift=10pt, draw=tugb, fill=white,text width=0.67\columnwidth, inner ysep=0.6mm,thick] (htarget) {\# Salary (int, \textbf{target feature}), min-max vals [ 100, 310 ], total-vals [ 5 ]  };


\node[textstyle, below=of rule.south west, anchor=north west, yshift=-5pt] (rule){\textbf{\#\#\# Categorical Features:}};
\node[textstyle, right=of rule, xshift=-324pt,draw=tugb, fill=tugb!10, inner ysep=0.8mm,thick, text width=0.3\columnwidth] (hcatalog) {Experience, Gender, Address, Zip};

% %%%%%%%%%%%%%%%%%%%%%%%%%%%%%%%%%    
\node [draw=lightgray,fill=white,rectangle,rounded corners,minimum width=0.73\columnwidth, text width=0.73\columnwidth, right=of sys_msg, anchor=north, xshift=700pt] (box) {
    % Begin of code
    \shortstack[l]{\\ \\ \\ \\ \\ \\ \\
    \linenumber{\ 1}\texttug{import} \textblack{pandas} \texttug{as} \textblack{pd} \\
    \linenumber{\ 2}\texttug{from} \textblack{sklearn.impute} \texttug{import} \textblack{SimpleImputer}\\
    \linenumber{\ 3}\texttug{from} \textblack{sklearn.preprocessing} \texttug{import} \textblack{OneHotEncoder}\\
    \linenumber{\ 4}\texttug{from} \textblack{sklearn.compose} \texttug{import} \textblack{ColumnTransformer}\\
    \linenumber{\ 5}\texttug{from} \textblack{sklearn.pipeline} \texttug{import} \textblack{Pipeline}\\
    \linenumber{\ 6}\texttug{from} \textblack{sklearn.ensemble} \texttug{import} \textblack{RandomForestRegressor}\\
    \linenumber{\ 7}\texttug{from} \textblack{sklearn.metrics} \texttug{import} \textblack{r2\_score} \\ \\ \\ \\ 
    \linenumber{\ 8}\textblack{train\_data = pd.read\_csv("}\textorang{train\_data.csv}\textblack{")} \\
    \linenumber{\ 9}\textblack{test\_data = pd.read\_csv("}\textorang{test\_data.csv}\textblack{")} \\ \\ \\ 
    \linenumber{10}\textblack{categorical\_features = ["\textorang{Experience}", "\textorang{Gender}", "\textorang{Address}", "\textorang{Zip}"]}\\
    \linenumber{11}\textblack{categorical\_transformer = Pipeline(steps=[}\\
    \textblack{\ \ \ \ \ \ \ \ \ \ \ \ \ \ ("\textorang{imputer}", SimpleImputer(strategy="\textorang{most\_frequent}")),} \\
    \textblack{\ \ \ \ \ \ \ \ \ \ \ \ \ \ ("\textorang{onehot}", OneHotEncoder(handle\_unknown="\textorang{ignore}"))])} \\ \\ 
    %
    \linenumber{12}\textblack{preprocessor = ColumnTransformer(}\\
    \textblack{\ \ \ \ \ \ \ \ \ \ \ \ \ \ transformers = [("\textorang{cat}", categorical\_transformer, categorical\_features)],} \\
    \textblack{\ \ \ \ \ \ \ \ \ \ \ \ \ \ remainder="\textorang{drop}" )} \\ \\
    \linenumber{13}\textblack{model = RandomForestRegressor(n\_estimators=\texttugb{100}, random\_state=\texttugb{42})} \\
    \linenumber{14}\textblack{pipeline = Pipeline(steps=[ ("\textorang{preprocessor}", preprocessor), ("\textorang{model}", model)])} \\ \\ \\
    %
    \linenumber{15}\textblack{X\_train = train\_data.drop("\textorang{Salary}", axis=\texttugb{1})} \\
    \linenumber{16}\textblack{y\_train = train\_data["\textorang{Salary}"]} \\
    \linenumber{17}\textblack{X\_test = test\_data.drop("\textorang{Salary}", axis=\texttugb{1})} \\
    \linenumber{18}\textblack{y\_test = test\_data["\textorang{Salary}"]} \\ \\ \\
    \linenumber{19}\textblack{pipeline.fit(X\_train, y\_train)} \\
    \linenumber{20}\textblack{y\_test\_pred = pipeline.predict(X\_test)} \\ \\ \\
    \linenumber{21}\textblack{test\_r2 = r2\_score(y\_test, y\_test\_pred)} \\
    \linenumber{22}\textblack{print(f"Test $R^2$: \{test\_r2\}")}
    }
}; % End of code
\node[circle,draw=lightgray!80,fill=lightgray!80, inner sep=0mm, minimum size=1.3em] (c1) at ($(box.north west)+(12pt,-12pt)$){};
\node[circle,draw=lightgray!80,fill=lightgray!80, inner sep=0mm, minimum size=1.3em, right=of c1, xshift=3pt] (c2) {};
\node[circle,draw=lightgray!80,fill=lightgray!80, inner sep=0mm, minimum size=1.3em, right=of c2, xshift=3pt] (c3) {};
\drawshadow{box}{lightgray}

%------------------------
\draw[-{Latex[scale=1]}, very thick, black] ($(ruleTask)+(38pt,0pt)$) -- ($(ruleTask)+(286pt,0pt)$);

\node[circle,draw=black,fill=black, inner sep=0mm, minimum size=1.7em, very thick, text=white, right=of hfe, xshift=97pt] (cclean)  {\textbf{\fontsize{17pt}{17pt}\selectfont{+}}};

\node[circle,draw=black,fill=black, inner sep=0mm, minimum size=1.7em,very thick, text=white] (cmodel) at($(box.north west)+(-33pt,-262pt)$) {\textbf{\fontsize{17pt}{17pt}\selectfont{+}}};

\draw [line width=2pt,decorate,decoration={calligraphic brace,amplitude=5pt}] ($(box.north west)+(3pt,-95pt)$) -- ($(box.north west)+(3pt,-24pt)$);

\draw [line width=2pt,decorate,decoration={calligraphic brace,amplitude=5pt}] ($(box.north west)+(3pt,-125pt)$) -- ($(box.north west)+(3pt,-105pt)$);

\draw [line width=2pt,decorate,decoration={calligraphic brace,amplitude=5pt}] ($(box.north west)+(3pt,-207pt)$) -- ($(box.north west)+(3pt,-132pt)$);

\draw [line width=2pt,decorate,decoration={calligraphic brace,amplitude=5pt}] ($(box.north west)+(3pt,-310pt)$) -- ($(box.north west)+(3pt,-213pt)$);

\draw [line width=2pt,decorate,decoration={calligraphic brace,amplitude=5pt}] ($(box.north west)+(3pt,-337pt)$) -- ($(box.north west)+(3pt,-317pt)$);

\draw[-{Latex[scale=1]}, very thick, black] (himport.east) --($(himport.east)+(292pt,0pt)$);

\draw[-{Latex[scale=1]}, very thick, black] (hdatasets.east) -| ($(hdatasets.east)+(100pt,-35pt)$) -- ($(hdatasets.east)+(130pt,-35pt)$);

\draw[-{Latex[scale=1]}, lightgreencode, very thick] (hclean.east) -| (cclean.north);
\draw[-{Latex[scale=1]}, tugb, very thick] ($(cclean.south)+(-3pt, -100pt)$) -- ($(cclean.south)+(-3pt, 0pt)$);
\draw[-{Latex[scale=1]}, lightgreencode, very thick] (hfe) -- (cclean.west);
\draw[-{Latex[scale=1]}, very thick, black] (cclean) -|($(cclean.east)+(30pt,-24pt)$) -- ($(cclean.east)+(53pt,-24pt)$);
\draw[-{Latex[scale=1]}, tugb, very thick] (hcatalog.east) -| ($(cclean.south)+(3pt, 0pt)$);

\draw[-{Latex[scale=1]}, lightgreencode, very thick] (hmodel.east) -| (cmodel.north);
\draw[-{Latex[scale=1]}, tugb, very thick] (htarget.east) -| (cmodel.south);
\draw[-{Latex[scale=1]}, very thick, black] (cmodel) -- ($(cmodel.east)+(22pt,0pt)$);

\draw[-{Latex[scale=1]}, very thick, black] (rule-ev.east) -| ($(rule-ev.west)+(455pt,-135pt)$) -- ($(rule-ev.west)+(501pt,-135pt)$);

\end{tikzpicture}